\paragraph{From a signed mean to every translate at the same exponent.}
The right side of \eqref{scope:missing} leaves more room per thick
coordinate when the thick set is made smaller. The original host conditions
allow a fixed proportion of that set to be retained. This permits an exact
transfer even when equality holds in \eqref{scope:average}.

\begin{proposition}[Reselection on the original windows]\label{scope:reselection}
Fix \(\psi,\varepsilon\) as in Section~\ref{ar:section}. Suppose that on a
set of scales the original common data satisfy \eqref{scope:average},
the full page rank bounds, the top cardinality bounds, the Behrend lower
bound with a uniform constant, and the hypotheses of
Proposition~\ref{ar:mean}. Retain the cyclotomic host and the construction
which, for each admissible choice of its two finite sets, supplies the
common centres and the points with their actual Harish--Chandra weights.
For assembly into a family, assume this construction is available on the
same set of scales with the new uniform Behrend constant. In the formal
construction one retains the input hosts of \texttt{HostExW}, the supply
of centres satisfying \texttt{WinOK} and the duty conditions, and
\texttt{ObjCP} for all the updated \texttt{HostGood} records before
\texttt{kitS}. The rank budget uses the field and rank of that host.

At every sufficiently late scale in this set one can choose new finite
sets and common points for which \eqref{scope:missing} holds for every
\(u\in A_0\), with its displayed coefficient. Both auxiliary primes,
the field, the ranks and \(\varepsilon\) at that scale are the original
ones. The new family also satisfies \eqref{scope:average}. If the
input gives arbitrarily late complete windows, the output has the
same complete windows after a uniform threshold.
\end{proposition}

\begin{proof}
Write \(t=|T|\), \(s=|T_2|\), \(y=x^\varepsilon\), \(p=y^{-1/8}\),
\(\delta=s/n\). Normalize a discrepancy or its budget by multiplying
by \(2/(tmb\delta)\). The normalized budget is exactly
\[
 W(t,s)=2y^5\left(\frac{A_1}{s}+\frac{A_3}{ts}+\frac{A_4}{t}\right),
 \qquad
 A_1=\frac{n a_E(1+m^2/k)}{mb},\quad
 A_3=2n^2my,\quad A_4=\frac{ky^2}{n}.
\]
The ranks satisfy \(ky^2\le n^2\sqrt{x}\) eventually, by
\(k\le2x^{9/2+10\varepsilon}\) and \(n\ge x^{2+8\varepsilon}\).
Also \(my\ge\sqrt{x}\) since \(m\ge\sqrt{xp}\). Thus
\(A_4s\le A_3\) for \(s\le n\). If a new choice has
\(t'\le4t\) and \(16s'\le s\), its first summand grows by a
factor of at least sixteen and its second by at least four.
The old third summand is bounded by the old second. Consequently
\[
                         W(t',s')\ge2W(t,s).
\]
Since \(t\le2np\), the second summand alone also gives
\(W(t,s)\ge 2\sqrt{x}\,y^{97/16}/\delta\).

Take any \(U\subset T_2\) of cardinality \(s'=\lfloor s/32\rfloor\)
in the original integer representatives. For \(s\ge64\) we have
\(s/64\le s'\le s/32\). The short range and the prohibition on
three-point relations in Proposition~\ref{host:behrend} pass to \(U\).
Put \(z=\sqrt{\log(n/o)\log L_w}\). If the old cardinality is at least
\((n/o)e^{-Cz}/L_w\), the new one is at least
\((n/o)e^{-(C+1)z}/L_w\) once \(z\ge\log64\). Also
\(\delta'=s'/n\ge x^{-\varepsilon}\) eventually, since
\(\log(\delta^{-1})=o(\log x)\) uniformly at these scales.
The upper bound \(4L_ws'\le n/o\) passes to subsets as well.

With \(U\) fixed, draw \(T'\) by Lemma~\ref{host:shares},
using probabilities \(p_v=p\) on \(A_0\) and \(p_v=3p/4\) elsewhere.
Write \(N(v)=|T'\cap(v+U)|\). These are binomial counts of mean
\(p_vs'\). For \(\lambda=\sqrt{12ps'\log(2n)}\), a Chernoff bound and a
union bound give \(\max_v|N(v)-p_vs'|\le\lambda\) with probability
\(1-o(1)\). The required coset shares and \(np/2\le t'\le2np\) also
hold with probability \(1-o(1)\). The original top cost exceeds its
prescribed bound with probability at most \(1/4\).
Choose a single outcome satisfying all of these conditions.
In particular \(t'\le4t\). The updated host has the same arithmetic
field, Artin data, level primes and thin-character bounds. In
\texttt{HostGood} the changed requirements are the two cardinalities,
the Behrend bound and the top cost. In \texttt{EAtH} only the coset
shares use \(T'\). Both \(E\) and its rank cutoff in
\(\mathrm{VVle}(E,\mathtt{R.t},\sqrt b)\) retain their values.

Write \(S=S_{\psi;q,d}(x)\) for the signed four-twist sum in
Proposition~\ref{ar:mean}. Set \(V=\exp(x^{\varepsilon/3})\),
\(\vartheta=\exp(-c x^{\varepsilon/8})\) with the fixed thin-bound
constant \(c>0\), and \(L=\log(2V)+n\vartheta\). In the additive
coordinates of Lemma~\ref{ar:balance}, put
\[
 K_h(z')=\frac{1}{2n}\Re\sum_{\chi}
                  \overline{\chi(1,0,h)}\,G_\chi(z').
\]
Subtract \(z'/V\) times each ramp at \(V\). For \(r=\ell^j\le V\)
the inverse Fourier coefficient is
\(Q_r(h)=2|I_\ell|^{-1}
\mathbf1_{(1,0,h)-jF_\ell\in I_\ell}\), including ramified primes.
It is nonnegative and \(\sum_h Q_r(h)\le2\). The resulting kernel is
\[
 K_h(z')=\frac{z'}{V}K_h(V)+\frac{z'\log(V/z')}{2n}
 -\sum_{r\le V}\Lambda(r)Q_r(h)\bigl(\min(1,z'/r)-z'/V\bigr).
\]
The thin bound gives \(|K_h(V)|\le2V\vartheta\).
Using \(\sum_{r\le V}\Lambda(r)/r\ll\log(2V)\) and
\(\sum_{r\le V}\Lambda(r)\ll V\), we obtain
\(\sum_h|K_h(x)|\ll xL\) and
\(\sum_h\sup_{|z'-x|\le w}|K_h(z')-K_h(x)|\ll wL\)
for \(w=\max_{t_0\in T'}|a_{t_0}/b-x|\le1\).

Use the deterministic convolution calculation in the proof of
Corollary~\ref{orb:overlap} with the present \(\lambda\). After
thick replacement and thin deletion the multiplicity at \(h\) is
\(N(u-h)\), with expectation \(s'p_{-h}\) for \(u\in A_0\).
The normalized new discrepancy \(d'_u\) satisfies simultaneously
\[
 |d'_u-S|\ll
      xL\frac{\lambda}{ps'}+wL+\frac1{x\delta'} .
\]
Here \(L\ll x^{\varepsilon/3}\) and \(ps'\gg x^{2+55\varepsilon/8}\).
The first term is \(O(x^{-149\varepsilon/48}\sqrt{\log(2x)})\).
Every term is \(o(W(t,s))\). This estimate uses only the genuine
weight windows. In the original construction
\(w_j(\tau)=B_\tau+(2j-1)\lceil x^{c_h}\rceil\) for \(1\le j\le m\).
The estimate allows the actual thick-place corrections at every point.

Write \(\bar d\) for the old normalized mean. Proposition~\ref{ar:mean}
gives \(|S-\bar d|\ll(1+x^\varepsilon)\log(2x)+1/(x\delta)
=o(W(t,s))\). The primes \(q,d\) in the two applications are identical.
By hypothesis \(\bar d\le W(t,s)\). Hence
\[
       d'_u\le W(t,s)+o(W(t,s))\le W(t',s')\qquad(u\in A_0)
\]
uniformly after a single threshold. Undo the new normalization to
obtain \eqref{scope:missing} and average it to obtain \eqref{scope:average}.

In the family version, the point-construction hypothesis is used after
\(U,T'\) have been fixed. The relation \texttt{Pt.From} carries the
actual top and thick sets, centres and raw field to every point.
The admissible host predicate can therefore retain the displayed
overlap condition on both sets. The construction
\texttt{exists\_famP\_of\_kits} preserves its input window set.
All translates at a given scale refer to the same new points and weights.
An isolated \texttt{RealizedData} record describes the old points and
contains no supply of new ones. Input estimates restricted to a proper
subinterval give the output on that subinterval. An independent window
theorem still requires a signed arithmetic estimate at every point
of a complete window.
\end{proof}

\begin{corollary}[An unbounded loss in the input mean]\label{scope:reselection-loss}
In Proposition~\ref{scope:reselection} multiply the old budget
in \eqref{scope:average} by \(g(x)\).
Suppose \(1\le g(x)\le \exp(a z)\) for some fixed \(a\ge0\)
and \(z=\sqrt{\log(n/o)\log L_w}\). The new pointwise and signed bounds
still have the original coefficient. In particular any fixed power
of \(\log x\) can be absorbed.
\end{corollary}
\begin{proof}
Take \(s'=\lfloor s/(32g(x))\rfloor\). Since \(s=n x^{-o(1)}\) and
\(g=x^{o(1)}\), we have \(s/(64g)\le s'\le s/(32g)\) eventually.
The Behrend lower bound survives with constant \(C+a+1\), and
\(s'/n\ge x^{-\varepsilon}\) still holds. In \(W(t',s')\) the first
two summands are at least \(32g\) and \(8g\) times the corresponding
old summands. Since the old third summand is bounded by the second,
\(W(t',s')\ge2gW(t,s)\). Both Fourier comparison errors remain
\(o(W(t,s))\), uniformly on the same scales.

Conversely, for any two legal choices on these same ranks,
\(t/t'\le4\). Comparing the three nonnegative summands separately
gives \(W(t',s')/W(t,s)\le4\max(1,s/s')\le4n/s'\).
The Behrend bound of the new family makes this at most
\(4oL_w\exp(C' z)\) for its fixed constant \(C'\).
Changing these sets can absorb losses with \(\log g=O(z)\).
A fixed positive power of \(x\) exceeds every budget amplification
of this form.
\end{proof}

If all large-order orbits must pass, Corollary~\ref{orb:conditional-overlap}
gives one choice with the same order of overlap error. For this choice
the original \texttt{BadOK} cost vanishes because it counts failed
large-order orbit tests. The budget comparison above applies directly.

The companion file \texttt{TranslateUniformity.lean} checks the subset
and Behrend inequalities, the updated host and character records, the
transport of both sets to the chosen points, and the budget comparison.
